% ============================================================================
% Chapter 2: Junction Symmetries
% ============================================================================
% Source: M6_GEOMETRY_DERIVATION.md, Z3_SYMMETRY_ANALYSIS_NEUTRON.md,
%         04c_routeB_z6_steiner.tex
% Status: FILLED from SoT
% Contamination check: PASSED (no SM terms)
% ============================================================================

\chapter{Junction Symmetries}
\label{ch:junction}

\source{M6\_GEOMETRY\_DERIVATION.md, Z3\_SYMMETRY\_ANALYSIS\_NEUTRON.md}

\EDCPolicyLockSentence

\begin{abstract}
How does the $\Zsix$ symmetry of the \AnchorJunction\ emerge from 5D topology?
This chapter formalizes the crystallization mechanism: boundary conditions on
$\Sfive$ select discrete symmetry groups via packing constraints. We derive
the subgroup chain $\Zsix \supset \Zthree \supset \Ztwo$ and interpret each
sector in EDC-native terms. The chapter closes by introducing the $\Msix$
coordination lattice---the dual graph of the Steiner junction network---which
bridges to the coordination and \Frustration\ framework of later chapters.
\end{abstract}

% ============================================================================
\section{Crystallization from $\Sfive$}
\label{sec:junction:crystallization}
% ============================================================================

\subsection{What ``Crystallization'' Means in EDC}

In standard condensed-matter physics, crystallization refers to the emergence
of discrete translational and rotational symmetries from continuous liquid
disorder. In EDC, the analogous process occurs in the configuration space of
brane defects:

\begin{derivationbox}[title=Definition: Junction Crystallization]
    \textbf{Junction crystallization} is the selection of a discrete symmetry
    group $G_{\text{jun}}$ from continuous 5D diffeomorphism invariance via:
    \begin{enumerate}
        \item Boundary conditions on the brane worldvolume
        \item Energy minimization (Nambu--Goto action)
        \item Topological constraints (non-trivial $\pione$)
    \end{enumerate}
    \tagDef
\end{derivationbox}

The outcome: Y-junction configurations self-organize into a discrete lattice
characterized by the symmetry group $G_{\text{jun}} \cong \Zsix$.

\subsection{Route B Recap: From P2 to $\Zsix$}

Chapter~\ref{ch:anchor} established the \AnchorJunction\ as a topological
ground state via two convergent routes. Route B (crystallographic) provides
the relevant framework here:

\begin{enumerate}
    \item \textbf{[P]} Postulate P2: Worldsheets have short-range repulsion
          and long-range confinement, with minimum at characteristic
          distance $d_0 > 0$.
    \item \textbf{[M]} Kepler--Hales theorem: Optimal 2D packing of equal
          circles is hexagonal, with packing fraction $\pi/(2\sqrt{3})$.
    \item \textbf{[Dc]} Crystallization: Worldsheet defects arrange in
          hexagonal lattice to minimize total energy.
    \item \textbf{[Dc]} $\Zthree \subset \Zsix$ at Y-junctions: The
          three-fold rotational symmetry of Y-vertices, combined with
          hexagonal lattice symmetry, yields $\Zsix$.
\end{enumerate}

\subsection{The Crystallization Claim}

\begin{resultbox}[title=Crystallization Result]
    The junction symmetry group is:
    \begin{equation}
        G_{\text{jun}} \cong \Zsix
        \label{eq:Gjun_Z6}
    \end{equation}
    This emerges from hexagonal packing of worldsheet defects on the
    $\Sfive$ boundary under Postulate P2. \tagDc
\end{resultbox}

\begin{edcopenbox}[title=OPEN: Explicit $\Sfive$ Derivation]
    Derive the explicit map from 5D action + thick-brane boundary conditions
    to $\Zsix$ crystallization. Current status: the result follows from P2 +
    Kepler--Hales, but a first-principles calculation from $\mathcal{E}[\Psi]$
    remains open. \tagOpen
\end{edcopenbox}

% ============================================================================
\section{Subgroup Structure}
\label{sec:junction:subgroups}
% ============================================================================

\subsection{Mathematical Statement}

The $\Zsix$ group has a well-known subgroup chain:

\begin{derivationbox}[title=Subgroup Chain]
    \begin{equation}
        \Zsix \supset \Zthree \supset \{e\}
        \quad\text{and}\quad
        \Zsix \supset \Ztwo \supset \{e\}
        \label{eq:subgroup_chain}
    \end{equation}
    Equivalently, $\Zsix \cong \Zthree \times \Ztwo$ (direct product for
    coprime orders 3 and 2). \tagDer
\end{derivationbox}

\begin{proof}
The group $\Zsix = \langle g \mid g^6 = e \rangle$ has order 6. By Lagrange's
theorem, subgroups have orders dividing 6: $\{1, 2, 3, 6\}$. The unique
subgroups are:
\begin{itemize}
    \item $\Zthree = \langle g^2 \rangle = \{e, g^2, g^4\}$ (index 2)
    \item $\Ztwo = \langle g^3 \rangle = \{e, g^3\}$ (index 3)
\end{itemize}
Since $\gcd(2,3) = 1$, the Chinese Remainder Theorem gives
$\Zsix \cong \Zthree \times \Ztwo$.
\end{proof}

\subsection{EDC-Native Interpretation}

Each subgroup has a distinct physical meaning in the junction framework:

\begin{table}[htbp]
\centering
\caption{Subgroup interpretation in junction space.}
\label{tab:subgroup_interpret}
\begin{tabular}{llll}
    \toprule
    \textbf{Group} & \textbf{Generator} & \textbf{Physical meaning} & \textbf{Tag} \\
    \midrule
    $\Zsix$ & $60°$ rotation & Full junction symmetry & [Dc] \\
    $\Zthree$ & $120°$ rotation & Arm permutation (cyclic) & [Dc] \\
    $\Ztwo$ & $180°$ rotation & Phase/parity sector & [P] \\
    \bottomrule
\end{tabular}
\end{table}

\subsubsection{$\Zthree$: Arm Permutation Symmetry}

The Y-junction has three arms emanating from the central hub. The $\Zthree$
subgroup acts by cyclic permutation:

\begin{derivationbox}[title=$\Zthree$ Action on Configurations]
    Let $\mathcal{C} = (L_1, L_2, L_3, \theta_{12}, \theta_{23}, \theta_{31})$
    denote a Y-junction configuration. The $\Zthree$ generator $g$ acts as:
    \begin{align}
        g \cdot (L_1, L_2, L_3) &= (L_2, L_3, L_1) \\
        g \cdot (\theta_{12}, \theta_{23}, \theta_{31}) &= (\theta_{23}, \theta_{31}, \theta_{12})
    \end{align}
    The group is $\Zthree = \{e, g, g^2\}$ with $g^3 = e$. \tagDc
\end{derivationbox}

\textbf{Physical consequence:} A configuration is $\Zthree$-symmetric if and
only if all three arms are equivalent: $L_1 = L_2 = L_3$ and
$\theta_{12} = \theta_{23} = \theta_{31} = 120°$.

\subsubsection{$\Ztwo$: Phase Sector}

The $\Ztwo$ factor corresponds to a discrete parity or phase flip. Two
interpretations are compatible with the mathematics:

\begin{enumerate}
    \item \textbf{Geometric:} $180°$ rotation about an axis perpendicular
          to the junction plane (exchanges arm pairs).
    \item \textbf{Topological:} Sign of the winding number or orientation
          of the worldsheet.
\end{enumerate}

\begin{edcopenbox}[title=OPEN: $\Ztwo$ Carrier Identification]
    The precise physical carrier of the $\Ztwo$ factor is not fully pinned
    down. Candidate interpretations include worldsheet orientation, boundary
    condition parity, or a discrete gauge sector. \tagOpen
\end{edcopenbox}

\subsection{Connection to Junction States}

The subgroup structure maps to junction configurations:

\begin{resultbox}[title=Junction--Subgroup Correspondence]
    \begin{itemize}
        \item \textbf{\AnchorJunction:} Full $\Zsix$ symmetry.
              Ground state with Steiner $120°$ geometry.
              (Chapter~\ref{ch:anchor})
        \item \textbf{\MetastableJunction:} $\Zthree$ symmetric branch.
              Excited state in double-well potential.
              (Chapter~\ref{ch:metastable})
        \item \textbf{$\Ztwo$ sector:} Distinguishes junction from anti-junction
              (if applicable).
    \end{itemize}
    \tagDc
\end{resultbox}

The key point: the \MetastableJunction\ is not a ``different particle'' but
a different \emph{configuration within the same junction space}, characterized
by which subgroup symmetry it preserves.

% ============================================================================
\section{$\Msix$ Lattice Structure}
\label{sec:junction:M6}
% ============================================================================

\subsection{From Junctions to Coordination}

When multiple junctions interact, they form networks. The language of
\emph{coordination}---how many neighbors each site has---provides the bridge
from single-junction symmetry to multi-junction cluster physics.

\begin{derivationbox}[title=Definition: $\Msix$ Coordination Lattice]
    The \textbf{$\Msix$ lattice} is the dual graph $G^*$ of the Steiner
    Y-junction network $G$:
    \begin{itemize}
        \item \textbf{Vertices of $G$:} Y-junctions (junction states)
        \item \textbf{Edges of $G$:} Flux tube bonds (leg-to-hub contacts)
        \item \textbf{Vertices of $G^*$:} Volumetric cells between junctions
        \item \textbf{Edges of $G^*$:} Contacts between adjacent cells
    \end{itemize}
    \tagDef
\end{derivationbox}

\subsection{Why $n = 6$: The Duality Argument}

The coordination number $n = 6$ in the $\Msix$ lattice is \emph{derived},
not postulated, from Steiner geometry:

\begin{derivationbox}[title=Derivation: $n = 6$ from Graph Duality]
    \textbf{Claim:} Each vertex in the dual graph $G^*$ has degree 6.

    \textbf{Proof:}
    \begin{enumerate}
        \item The primal graph $G$ is 3-valent (each Y-junction has 3 legs).
        \item At each vertex of $G$, the three edges meet at $120°$ angles.
        \item The exterior angle at each vertex is $180° - 120° = 60°$.
        \item For a closed face, total exterior angle $= 360°$.
        \item Minimum vertices per face $= 360°/60° = 6$.
        \item Therefore, minimal faces of $G$ are \textbf{hexagons} (6-cycles).
        \item By graph duality: degree of vertex in $G^*$ = edges of
              corresponding face in $G$ = 6. \qed
    \end{enumerate}
    \tagDer
    \label{deriv:n_equals_6}
\end{derivationbox}

\begin{figure}[htbp]
\centering
\begin{verbatim}
      Primal G (3-valent)          Dual G* (6-valent)

         J1 --- J2                     *---*---*
        / |     | \                   / \ / \ / \
      J6  |     |  J3                *---*---*---*
        \ |     | /                   \ / \ / \ /
         J5 --- J4                     *---*---*

      Hexagonal faces             Each vertex has 6 neighbors
\end{verbatim}
\caption{Schematic of primal Steiner lattice $G$ and dual $\Msix$ lattice $G^*$.}
\label{fig:primal_dual}
\end{figure}

\subsection{Sites, Edges, and Pinning}

In the $\Msix$ language:

\begin{table}[htbp]
\centering
\caption{$\Msix$ lattice vocabulary.}
\label{tab:M6_vocab}
\begin{tabular}{lll}
    \toprule
    \textbf{$\Msix$ term} & \textbf{Definition} & \textbf{Physical role} \\
    \midrule
    Site & Vertex of $G^*$ (cell in primal) & Location for \Pinning \\
    Edge & Contact between adjacent cells & \Pinning\ interface \\
    Coordination $\ncoord$ & Number of edges at a site & Cluster characterization \\
    \bottomrule
\end{tabular}
\end{table}

The \textbf{\Pinning} energy depends on how many contacts a cluster has with
its environment. This motivates the coordination function $\ncoord(A)$
developed in Chapter~\ref{ch:M6}.

\subsection{Allowed Coordination Set}

Not all integers are achievable as coordination numbers. The $\Msix$ topology
constrains the allowed values:

\begin{derivationbox}[title=Allowed Coordination Set]
    The allowed coordination numbers form the set:
    \begin{equation}
        S = \{2^a \times 3^b \mid a, b \geq 0\}
        = \{1, 2, 3, 4, 6, 8, 9, 12, 16, 18, 24, 27, 32, 36, \ldots\}
        \label{eq:allowed_set}
    \end{equation}
    Integers not in $S$ correspond to \textbf{frustrated} configurations.
    \tagDer
\end{derivationbox}

\textbf{Forward reference:} The derivation of $S$ from $\Msix$ topology and
its physical consequences (frustration, forbidden zones) are developed in
Chapters~\ref{ch:M6} and~\ref{ch:frustrated}.

\subsection{The Forbidden Zone}

A key prediction of the $\Msix$ framework is the existence of gaps in the
allowed coordination spectrum:

\begin{resultbox}[title=Forbidden Zone Preview]
    The \ForbiddenZone\ $[37, 47]$ contains 11 consecutive integers, none
    expressible as $2^a \times 3^b$. Clusters with coordination in this
    range experience maximal \Frustration.
    \tagDer\ (derived in Chapter~\ref{ch:frustrated})
\end{resultbox}

% ============================================================================
\section{Summary}
\label{sec:junction:summary}
% ============================================================================

\begin{resultbox}[title=Chapter 2 Summary: Junction Symmetries]
    \textbf{Derivation chain:}
    \[
        \Sfive \text{ boundary} \xrightarrow{\text{P2 + packing}}
        \text{crystallization} \xrightarrow{\text{hex lattice}}
        \Zsix \text{ symmetry} \xrightarrow{\text{subgroups}}
        \Zthree \text{ branch} \xrightarrow{\text{dual graph}}
        \Msix \text{ coordination}
    \]

    \textbf{Key results:}
    \begin{itemize}
        \item $G_{\text{jun}} \cong \Zsix$ from hexagonal crystallization [Dc]
        \item $\Zsix \cong \Zthree \times \Ztwo$ with distinct physical roles [Der]
        \item \AnchorJunction\ = full $\Zsix$; \MetastableJunction\ = $\Zthree$ branch [Dc]
        \item $\Msix$ = dual of Steiner lattice; coordination $n = 6$ [Der]
        \item Allowed set $S = \{2^a \times 3^b\}$ [Der]
    \end{itemize}

    \textbf{Forward connections:}
    \begin{itemize}
        \item Chapter~\ref{ch:metastable}: $\Zthree$ branch and double-well dynamics
        \item Chapter~\ref{ch:M6}: Full $\Msix$ coordination derivation
        \item Chapter~\ref{ch:frustrated}: \Frustration\ correction from $d(\ncoord)$
    \end{itemize}
\end{resultbox}

% ============================================================================
\section{Epistemic Status}
\label{sec:junction:epistemic}
% ============================================================================

\begin{table}[htbp]
\centering
\caption{Epistemic status of Chapter 2 claims.}
\label{tab:ch2_epistemic}
\begin{tabular}{llll}
    \toprule
    \textbf{Claim} & \textbf{Status} & \textbf{Reference} & \textbf{Source} \\
    \midrule
    Crystallization definition & [Def] & \S\ref{sec:junction:crystallization} & — \\
    $G_{\text{jun}} \cong \Zsix$ & [Dc] & Eq.~\eqref{eq:Gjun_Z6} & P2 + Kepler--Hales \\
    Subgroup chain $\Zsix \supset \Zthree$ & [Der] & Eq.~\eqref{eq:subgroup_chain} & Group theory \\
    $\Zthree$ = arm permutation & [Dc] & \S\ref{sec:junction:subgroups} & Configuration analysis \\
    $\Ztwo$ carrier identity & [P] & \S\ref{sec:junction:subgroups} & Interpretation \\
    \midrule
    $\Msix$ lattice definition & [Def] & \S\ref{sec:junction:M6} & — \\
    $n = 6$ from duality & [Der] & Derivation~\ref{deriv:n_equals_6} & Graph duality \\
    Allowed set $S = \{2^a 3^b\}$ & [Der] & Eq.~\eqref{eq:allowed_set} & $\Msix$ topology \\
    \midrule
    Explicit $\Sfive \to \Zsix$ map & [OPEN] & \S\ref{sec:junction:crystallization} & Future work \\
    $\Ztwo$ physical carrier & [OPEN] & \S\ref{sec:junction:subgroups} & Future work \\
    \bottomrule
\end{tabular}
\end{table}

% ============================================================================
\section{Notation Summary}
\label{sec:junction:notation}
% ============================================================================

\begin{table}[htbp]
\centering
\caption{Chapter 2 notation dictionary.}
\label{tab:ch2_notation}
\begin{tabular}{lll}
    \toprule
    \textbf{Symbol} & \textbf{Meaning} & \textbf{Status} \\
    \midrule
    $G_{\text{jun}}$ & Junction symmetry group & [Dc] \\
    $\Zsix$ & Hexagonal symmetry; $\cong \mathbb{Z}/6\mathbb{Z}$ & [Der] \\
    $\Zthree$ & Arm permutation subgroup; $\cong \mathbb{Z}/3\mathbb{Z}$ & [Der] \\
    $\Ztwo$ & Phase/parity subgroup; $\cong \mathbb{Z}/2\mathbb{Z}$ & [P] \\
    $G$ & Primal Steiner graph (3-valent) & [Def] \\
    $G^*$ & Dual graph = $\Msix$ lattice (6-valent) & [Def] \\
    $\Msix$ & Coordination lattice & [Def] \\
    $\ncoord$ & Coordination number & [Def] \\
    $S$ & Allowed coordination set $\{2^a 3^b\}$ & [Der] \\
    \bottomrule
\end{tabular}
\end{table}

